\documentclass[11pt,a4paper]{article}

\usepackage[T1]{fontenc}
\usepackage{lmodern}
\usepackage[utf8]{inputenc}
\usepackage{textcomp}
\usepackage{amsmath,amssymb,amsthm}
\usepackage{mathtools}
\usepackage{hyperref}
\usepackage{doi}
\usepackage{geometry}
\usepackage{booktabs}
\usepackage{array}
\usepackage{enumitem}

\input{glyphtounicode}
\pdfgentounicode=1

\geometry{margin=2.5cm}

\hypersetup{
  colorlinks=true,
  linkcolor=black,
  citecolor=black,
  urlcolor=black
}

\theoremstyle{definition}
\newtheorem{remark}{Remark}

\newcommand{\SU}{\mathrm{SU}}
\newcommand{\UU}{\mathrm{U}}
\newcommand{\Epi}{E_\Pi}
\newcommand{\Sbundle}{\mathcal{S}_\Pi}
\newcommand{\Upmns}{U_{\mathrm{PMNS}}}
\newcommand{\ket}[1]{|#1\rangle}
\newcommand{\Uckm}{U_{\mathrm{CKM}}}

\title{Neutrino Mixing and Projective Particle Identity\\[0.4em]
\large A Structural Note}

\author{J\'{e}r\^{o}me Beau\\[0.3em]
\small Independent Researcher\\
\small \texttt{jerome.beau@cosmochrony.org}}

\date{2026}

\begin{document}

  \maketitle

  \begin{abstract}
    Neutrino oscillations establish that the weak flavour basis and the mass propagation basis of the lepton
    sector are not aligned.
    This note examines what this misalignment means for the ontology of particle identity within the Cosmochrony
    framework, where particles are invariants of a non-injective projection rather than primitive objects.
    Three structural observations are developed.
    First, basis misalignment is generic: neutral meson systems exhibit the same distinction between interaction
    labels and propagating eigenmodes, and the neutrino is distinguished only by being elementary within the
    Standard Model inventory.
    Second, within the projective reading, alignment --- not misalignment --- is the feature requiring a
    mechanism: an identity-stabilising projection channel.
    Charged fermions possess such channels; neutrinos, singlets under $\SU(3)_C \times \UU(1)_{\mathrm{em}}$,
    do not.
    The large PMNS mixing is then the unfixed default, consistent with the anarchy hypothesis, while the
    near-diagonality of the CKM matrix becomes the non-trivial explanandum, requiring the charged channels to pin
    the up-type and down-type mass bases to a common projective frame.
    Third, the Dirac-or-Majorana alternative is reread as the degree to which the projection resolves relational
    conjugation, with neutrinoless double beta decay as its standard discriminator.
    The note is structural throughout: it derives no mixing angles and no masses.
    It fixes the reading under which the neutrino sector is the sharpest available test case for projective
    particle identity, and states one falsifiable structural expectation: sectors with fewer identity-stabilising
    projection channels generically exhibit larger mixing between interaction and propagation bases.
    Interpretive status: the reinterpretation of particle identity as a projection invariant is a structural
    reading of standard neutrino phenomenology, not a competing prediction for it; all quantitative content of
    the standard oscillation formalism is retained as the operational effective description.
  \end{abstract}

  \medskip
  \noindent\textbf{Keywords.}
  neutrino oscillations; particle identity; projective invariants; PMNS matrix; CKM matrix; anarchy hypothesis;
  neutral meson mixing; Dirac and Majorana neutrinos; non-injective projection; ontology of particle physics.

  \newpage
  \tableofcontents

  \newpage
  \section{Introduction}
    \label{sec:introduction}

    The Standard Model tabulates particles as primitive entries: each species is a well-defined object carrying a
    fixed set of labels.
    Neutrinos are the elementary case where this tabulation is under visible strain.
    The states produced and detected by the weak interaction ($\nu_e$, $\nu_\mu$, $\nu_\tau$) are not the states
    that propagate; the observed identity of a neutrino depends on which basis the observation accesses.
    This tension has recently been articulated as a pressure point on Standard-Model ontology in the philosophy
    of particle physics~\cite{Hobart2026OnticTypes}.

    The Cosmochrony framework~\cite{Cosmochrony} treats particles as invariants of a non-injective projection
    $\Pi$ from a static relational substrate, rather than as primitive objects.
    The white paper already assigns neutrinos a special position, as \emph{partially projectable modes} close to
    the boundary of projectability.
    The present note develops the structural core of that position in the language of the fermionic-matter
    sub-programme~\cite{Beau2026fmnote}, where the fermionic sector arises from the admissible spinor bundle
    $\Sbundle$ and its projective endomorphism $\Epi$~\cite{Beau2026q14}.

    The scope is deliberately narrow.
    The note does not derive mixing angles, masses, or oscillation corrections; the chiral polar factor carrying
    generation mixing is invisible to $\Epi^2$ at the present stratum~\cite{Beau2026pyo}, so any quantitative
    claim would be premature.
    What the note fixes is the reading: which features of the neutrino sector are generic, which require a
    mechanism, and which structural expectation is falsifiable.

  \section{The Standard Misalignment}
    \label{sec:standard-misalignment}

    The following facts are standard and are imported, not derived.
    Flavour eigenstates are defined by the charged-current weak interaction; mass eigenstates diagonalize
    propagation; the two bases are related by the unitary PMNS matrix $\Upmns$,
    \[
      \ket{\nu_\alpha} = \sum_i (\Upmns^*)_{\alpha i}\,\ket{\nu_i}, \qquad \alpha \in \{e,\mu,\tau\},
    \]
    a structure anticipated by Pontecorvo~\cite{Pontecorvo1957} and by Maki, Nakagawa and
    Sakata~\cite{MakiNakagawaSakata1962}, and established experimentally by the observation of atmospheric and
    solar neutrino oscillations~\cite{Fukuda1998, Ahmad2002}.
    Oscillations imply non-degenerate masses and therefore force an extension of the minimal Standard Model, in
    which neutrinos are massless.
    In the quark sector the analogous misalignment between up-type and down-type mass bases is encoded by the
    CKM matrix~\cite{Cabibbo1963, KobayashiMaskawa1973}, which is experimentally near-diagonal, whereas
    $\Upmns$ contains two large angles.

    Operationally, the standard oscillation formalism --- interference between mass eigenstates with distinct
    dispersion --- is retained without modification throughout this note.
    The question addressed here is not \emph{how} neutrinos oscillate but \emph{what oscillation says about the
    ontology of the oscillating object}.

  \section{The Neutral-Meson Comparison}
    \label{sec:neutral-mesons}

    Basis misalignment is not exceptional in itself.
    Neutral meson systems exhibit the same structural distinction: $K^0$ and $\bar K^0$ are the strong-interaction
    (production) labels, while the propagating eigenmodes are their mixtures, as recognised by Gell-Mann and
    Pais~\cite{GellMannPais1955}.
    Whenever an interaction basis and a propagation basis coexist and no mechanism forces them to coincide, the
    generic situation is that they do not.

    The ontological difference is that neutral mesons are composite: their mixing can be attributed to the
    dynamics of constituents, and the naive particle ontology can retreat to the quark level.
    Neutrinos are elementary within the Standard Model inventory, so no such retreat is available.
    The neutrino is therefore the case where the contextuality of particle identity --- flavour at
    production and detection, mass in propagation --- must be confronted at the level of the fundamental
    tabulation itself.

  \section{The Projective Reading: Alignment Requires a Stabiliser}
    \label{sec:projective-reading}

    In the Cosmochrony framework a particle label is an invariant of the projection $\Pi$: flavour is an
    invariant of the weak coupling, and mass is read as a spectral invariant of the projected residue, the eigenvalue
    data of the projective endomorphism $\Epi$~\cite{Beau2026q14}, through a generation reading map that the
    Schur reduction of the sub-programme takes as a hypothesis~\cite{Beau2026prs}.
    These are distinct invariants of distinct structures, and there is no structural reason for the weak flavour
    basis to diagonalize the projected spectral mass residue.
    The structural statement of this note is therefore:
    \begin{quote}
      \emph{Basis misalignment is the generic situation; basis alignment requires an additional projective
      stabiliser.}
    \end{quote}
    Neutrino oscillations are then not an anomaly of particle identity but the expected manifestation of a
    neutral chiral matching mode whose weak-interaction label and spectral mass label are distinct projective
    invariants.

    This restates, in the vocabulary of the fermionic-matter sub-programme, the white-paper description of
    neutrinos as partially projectable modes~\cite{Cosmochrony}.
    One white-paper formulation is deliberately softened here.
    The claim that oscillations reflect ``variations in projectability rather than interference between
    fundamental mass eigenstates'' is retained only in its ontological part: the standard interference
    description remains the operational effective description, and what the framework reinterprets is the status
    of the mass eigenstates --- projected spectral invariants rather than fundamental degrees of freedom.

    \paragraph{Relation to the white-paper description.}
    The white paper's neutrino subsection~\cite{Cosmochrony} contains, beyond the elements developed here,
    several further structural claims: the smallness of neutrino masses as a necessity of boundary proximity,
    the failure of absolute localization, the role of neutrino-like excitations as non-local channels of
    projective reorganization in decay processes, and their cosmological smoothing role.
    These are distinct claims, not required by the identity argument, and they are deliberately left outside the
    scope of this note.
    What is developed here is the identity-theoretic core only: partial projectability read as the absence of
    stabilising channels, the flavour and mass bases as distinct projective invariants, and the conjugation
    reading of the Dirac--Majorana alternative.

  \section{The CKM--PMNS Contrast}
    \label{sec:ckm-pmns}

    The projective reading reverses the usual explanatory burden.
    As a structural reading, not a derived result, charged fermions possess projection channels that stabilise
    their identity between production and detection --- the electromagnetic and colour couplings, together with
    the operational pinning provided by large mass splittings and rapid decoherence --- whereas neutrinos,
    singlets under $\SU(3)_C \times \UU(1)_{\mathrm{em}}$, lack any such stabilising channel.
    This is not a Standard-Model dynamical statement: electromagnetic and colour interactions do not by
    themselves diagonalize flavour.
    It is a projective reading of why charged sectors may acquire operationally stable identity labels.

    The large leptonic mixing encoded by $\Upmns$ is then the unfixed default.
    This is consistent with the anarchy hypothesis for neutrino masses~\cite{HallMurayamaWeiner2000}, which
    showed that the oscillation data are well accounted for by a mass matrix with no preferred structure.
    The projective reading adds to anarchy what anarchy does not supply: a reason for its sector selectivity.
    Randomness is expected precisely where no identity-stabilising channel constrains the basis, and only there.

    Conversely, the near-diagonality of $\Uckm$ becomes the non-trivial feature to be explained.
    In this reading, CKM near-alignment requires the additional assumption that the charged projection channels
    pin the up-type and down-type mass bases to a \emph{common} projective frame.
    This common-frame assumption is the load-bearing hypothesis of the reversal: two sectors stabilised
    independently could remain mutually rotated, so stabilisation alone does not imply alignment.
    Its status is conditional, and its derivation is an open problem of the fermionic-matter
    sub-programme~\cite{Beau2026fmnote}.

    The reversal yields one falsifiable structural expectation:
    \begin{quote}
      \emph{Sectors with fewer identity-stabilising projection channels should generically exhibit larger mixing
      between interaction and propagation bases, unless an independent alignment mechanism is present.}
    \end{quote}
    The known pattern --- large PMNS angles in the unpinned neutral sector, small CKM angles in the doubly
    pinned quark sector --- conforms to this expectation; a future extension of the inventory (for example a
    sterile or otherwise unpinned state) would test it further.

  \section{Dirac versus Majorana as Conjugation Resolution}
    \label{sec:dirac-majorana}

    The Majorana alternative~\cite{Majorana1937} asks whether the neutrino is distinct from its antiparticle.
    In the projective reading this is not a fundamental binary choice but a question about the resolution of the
    projection: a Dirac neutrino corresponds to a configuration whose conjugate projected realisations are
    distinguished by $\Pi$, a Majorana neutrino to one whose partial projectability collapses the conjugate
    classes~\cite{Cosmochrony}.
    Lepton number conservation is then an emergent bookkeeping of regimes where conjugate configurations remain
    distinguishable, and its violation becomes admissible exactly where the distinction is erased.

    The standard discriminator, neutrinoless double beta decay, acquires a structural interpretation: it probes
    whether the projection resolves relational conjugation in the neutral chiral sector.
    No rate prediction is made here; the observation or exclusion of the process constrains the resolution
    structure of $\Pi$, not a parameter of this note.

  \section{Phenomenological Outlook (Structural)}
    \label{sec:phenomenology}

    The following directions are recorded as structural expectations, without quantitative content at this
    stage.
    Long-baseline coherence: as a mode without identity-stabilising channels, the neutrino is the natural probe
    of decoherence effects sensitive to global properties of the projected geometry; departures from standard
    oscillation patterns at ultra-long baselines would be a possible corresponding signature.
    Neutrinoless double beta decay: the conjugation-resolution reading of Section~\ref{sec:dirac-majorana}.
    Cosmological backgrounds: as the modes closest to the projection boundary, relic neutrinos are the
    observational window most sensitive to pre-geometric structure~\cite{Cosmochrony}.
    Deriving any of these quantitatively requires the explicit Lorentzian $\Epi$ and the Yukawa sector, which
    are the primary open deliverables of the fermionic-matter sub-programme~\cite{Beau2026fmnote}.

  \section{Positioning}
    \label{sec:positioning}

    Three adjacent positions delimit the present one.
    Hobart's ontic-types analysis~\cite{Hobart2026OnticTypes} identifies the same pressure point --- neutrino
    mixing challenges a one-to-one correspondence between particle type, flavour label, and mass eigenstate ---
    and proposes a family ontology in which the manifest state depends on interaction conditions.
    The present note shares the contextuality diagnosis but supplies a mechanism candidate: identity is a
    projection invariant, and its contextuality is the absence of a stabilising channel, embedded in a framework
    where the fermionic sector itself is derived~\cite{Beau2026q14}.
    The anarchy hypothesis~\cite{HallMurayamaWeiner2000} shares the ``large mixing is the default'' conclusion;
    the present reading adds the sector-selectivity criterion and the reversal of the explanatory burden onto
    CKM.
    Standard quantum field theory treats flavour and mass bases as effective bookkeeping related by unitary
    rotation; the present note retains that formalism operationally and differs only in the ontological status
    assigned to the bases.

  \section{Status of Claims}
    \label{sec:status}

    Standard (imported): flavour--mass basis misalignment, oscillations, neutral meson mixing, near-diagonal
    CKM, large-angle PMNS, anarchy consistency of the oscillation data.
    Structural (this note): the stabiliser reading of alignment, the identification of the neutrino as a neutral
    chiral matching mode, the PMNS-default/CKM-explanandum reversal, the conjugation-resolution reading of the
    Dirac--Majorana alternative, and the mixing-versus-channels expectation.
    Conditional: the common projective frame assumption for CKM near-alignment (load-bearing, underived).
    Open: the quantitative pinning mechanism, the derivation of mixing magnitudes (blocked at the trivial chiral
    polar class $[U_\Pi]=[I]$ of the present stratum~\cite{Beau2026pyo}), the resolution criterion for the
    Dirac--Majorana alternative, and any oscillation correction.

  \section{Conclusion}
    \label{sec:conclusion}

    Neutrinos are the elementary case where particle identity is manifestly contextual: the weak basis, the mass
    basis, and the ontological inventory do not coincide by necessity.
    Read projectively, this contextuality is not a puzzle but the default: identity labels are invariants of a
    non-injective projection, distinct invariants need not align, and alignment is what requires a mechanism.
    The neutrino sector thereby becomes the sharpest available test case for projective particle identity, and
    the quark sector inherits the substantive question: why is CKM nearly diagonal?

    \subsection*{Interpretive outlook}
    The following is an interpretation, a structural reading of established phenomenology, and not a result.
    If particle identity is projective rather than primitive, the Standard Model table is a catalogue of stable
    projection invariants, and the neutrino marks its frontier: the entry least pinned by projection channels,
    closest to the boundary where the projected description itself degrades.
    On this reading, oscillation experiments are not merely measuring parameters of an extended Standard Model;
    they are probing how much identity a mode can retain when almost nothing in the projection holds it fixed.
    The interpretive wager of the framework is that the pattern of mixing across sectors --- anarchic where
    unpinned, aligned where pinned --- is a visible trace of the projection structure from which particles
    emerge.

  \bibliographystyle{plain}
  \bibliography{cosmochrony-bibliography,external-refs}

\end{document}
